\documentclass[11pt]{article}

\usepackage[a4paper,margin=1in]{geometry}
\usepackage{amsmath,amssymb,amsthm,mathtools,bm}
\usepackage{microtype}
\usepackage{enumitem}
\usepackage[colorlinks=true,linkcolor=blue,citecolor=blue,urlcolor=blue]{hyperref}
\usepackage[nameinlink,capitalize]{cleveref}

\newtheorem{theorem}{Theorem}[section]
\newtheorem{lemma}[theorem]{Lemma}
\newtheorem{proposition}[theorem]{Proposition}
\newtheorem{corollary}[theorem]{Corollary}
\theoremstyle{definition}
\newtheorem{definition}[theorem]{Definition}
\newtheorem{remark}[theorem]{Remark}

\newcommand{\R}{\mathbb{R}}
\newcommand{\Pcal}{\mathcal{P}}
\newcommand{\supp}{\operatorname{supp}}
\newcommand{\dd}{\,\mathrm{d}}
\newcommand{\ee}{\mathrm{e}}
\newcommand{\betaHPS}{\beta_{3}}
\newcommand{\kernel}{k}

\title{An Exact Evaluation of a Radial Mean--Field Constant\\in the Atomic Excess--Charge Problem}
\author{Hu Suiye}
\date{September 2026}

\begin{document}
\maketitle

\begin{abstract}
Hundertmark, Pattakos and Schulz introduced a mean--field constant $\beta_s$ in their recent refinement of the excess--charge bound for atoms.  For $1\le s\le 3$ they reduced the variational problem to radial probability measures, while for $s=3$ their explicit estimate used a pointwise lower bound that yields the coefficient $1.1185$ in the resulting ionization bound.  We solve the remaining one--dimensional radial variational problem at $s=3$ exactly.  After passing to logarithmic radius and an exponentially tilted measure, the minimization becomes a maximization problem for the translation--invariant kernel
\[
 k(t)=\ee^{-|t|}-\ee^{-2|t|}.
\]
The unique normalized maximizer is supported on a finite interval and has cosine density.  Writing
\[
 a_*:=\frac{1}{\sqrt2}\arctan\sqrt2,
\]
we obtain
\[
 \boxed{\displaystyle
 \beta_3=\frac{3\ee^{2a_*}}{1+3\ee^{2a_*}}
 =0.9205346822904167\ldots}
\]
and hence $\beta_3^{-1}=1.086325175181735\ldots$.  The proof uses an exact obstacle condition, compactness of maximizing sequences, endpoint smooth pasting, and a conditional negative--definiteness argument on the extremal support based on a Poincar\'e estimate.  We also record the implication for the leading constant in the Hundertmark--Pattakos--Schulz framework and isolate the remaining finite--$N$ bookkeeping needed for a sharpened atomic bound.
\end{abstract}

\section{Introduction}

The excess--charge problem asks how many electrons an atomic nucleus of charge $Z$ can bind.  In the nonrelativistic Schr\"odinger model one expects only a bounded number of excess electrons, independently of $Z$.  A classical theorem of Lieb gives $N_c(Z)<2Z+1$ \cite{Lieb1984}; Nam later improved the coefficient to $1.22$ with an explicit $O(Z^{1/3})$ correction \cite{Nam2012}.  The conjectural bound remains $N_c(Z)\le Z+C$ for a universal constant $C$; see, for example, \cite{Nam2013Review}.

A recent refinement by Hundertmark, Pattakos and Schulz (HPS) develops the weighted multiplier method substantially further and proves
\[
 N_c(Z)<1.1185\,Z+O(Z^{1/3}),
\]
with an explicit lower--order term \cite{HPS2025}.  A central object in their argument is the mean--field constant
\begin{equation}\label{eq:def-beta-s}
 \beta_s:=\inf_{\mu}
 \frac{\displaystyle
 \iint_{\R^3\times\R^3}
 \frac{|x|^s+|y|^s}{2|x-y|}\,\dd\mu(x)\dd\mu(y)}
 {\displaystyle\int_{\R^3}|x|^{s-1}\,\dd\mu(x)},
\end{equation}
where the infimum is taken over the admissible probability measures used in \cite{HPS2025}.  HPS prove that for $1\le s\le3$ the infimum can be restricted to radial measures.  For $s=3$ they then use a pointwise estimate on the radial kernel.  This yields the explicit lower bound
\[
 \beta_3\ge 0.8941074569\ldots,
\]
whose reciprocal is the source of the leading coefficient $1.1185$ in their numerical theorem.

The first purpose of the present paper is to solve the radial variational problem at $s=3$ exactly.

\begin{theorem}[Exact value of $\beta_3$]\label{thm:main}
Let $\beta_3$ be the HPS mean--field constant defined by \eqref{eq:def-beta-s}.  Set
\[
 a_*:=\frac{1}{\sqrt2}\arctan\sqrt2.
\]
Then
\begin{equation}\label{eq:beta-exact}
 \boxed{\displaystyle
 \beta_3=\frac{3\ee^{2a_*}}{1+3\ee^{2a_*}}}
 =0.9205346822904167\ldots .
\end{equation}
In particular,
\[
 \beta_3^{-1}=1.086325175181735\ldots .
\]
A minimizing radial probability measure is supported on the shell
\[
 \ee^{-a_*}\le |x|\le \ee^{a_*}
\]
and has spatial density proportional to
\[
 |x|^{-4}\cos\!\bigl(\sqrt2\log|x|\bigr).
\]
Among radial admissible measures, the minimizer is unique up to the scaling invariance of \eqref{eq:def-beta-s}.  No uniqueness assertion is made here for minimizers in the full nonradial HPS class.
\end{theorem}

The second purpose is to propagate \Cref{thm:main} through the finite--$N$ part of the HPS argument.  Their cubic-weight estimates retain the actual mean--field constant $\beta_3$ symbolically; the pointwise lower bound enters only when the final coefficients are converted to explicit numbers.  Thus the present finite--$Z$ improvement is not obtained by replacing a decimal in the final HPS theorem: it comes from inserting the exact value \eqref{eq:beta-exact} into their symbolic chain and redoing the endpoint optimization.

\begin{theorem}[Sharpened finite--$Z$ excess--charge bound]\label{thm:finite-Z}
For every $Z\ge4$,
\begin{equation}\label{eq:finite-Z-main}
\boxed{\begin{aligned}
N_c(Z)<{}&\beta_3^{-1}Z+3.812\,Z^{1/3}+0.01294\\
&+0.18185\,Z^{-1/3}+0.01941\,Z^{-2/3},
\end{aligned}}
\end{equation}
where $\beta_3$ is the exact constant in \eqref{eq:beta-exact}.  Consequently,
\begin{equation}\label{eq:finite-Z-simple}
 \boxed{N_c(Z)<1.086326\,Z+3.90\,Z^{1/3}},\qquad Z\ge4.
\end{equation}
\end{theorem}

\subsection*{Significance of the result}

The improvement has a direct quantitative meaning.  The HPS leading coefficient $1.1185$ is produced by a lower estimate on the cubic mean--field constant, whereas the present computation replaces that estimate by the exact radial value.  The leading coefficient therefore drops to
\[
 1.086325175\ldots .
\]
Relative to the conjectural coefficient $1$, the distance from the HPS value is $0.1185$, while the exact radial computation removes
\[
 1.1185-1.086325175\ldots=0.032174824\ldots,
\]
or about $27\%$ of that leading-order gap.  This comparison is only about the linear coefficient; the $Z^{1/3}$ correction remains part of the finite--$Z$ theorem.  The point is that this gain is obtained without a new many--body inequality.  It isolates the amount of the HPS loss that came specifically from replacing a global mean--field optimization by a pointwise pairwise bound.

There is also a structural meaning.  The pointwise HPS estimate asks for the worst ratio associated with one pair of radii, but a probability measure cannot make every pair simultaneously realize that worst geometry.  Solving the full variational problem enforces this global compatibility.  The optimizer is not a singular pair configuration: it is an extended density supported on a finite annulus, with the explicit log--radial cosine profile in \Cref{thm:main}.  Thus the improvement is a genuinely collective effect inside the scalar radial mean--field problem.

Finally, the extremizer identifies the obstruction that remains after the scalar radial problem is solved exactly.  In physical radius the density is a compactly supported cosine modulation of $r^{-4}$, the critical power already suggested by the cubic weighted multiplier structure.  Consequently, any further reduction of the leading coefficient within a strengthened cubic-weight argument must use information not present in the scalar radial probability measure alone---for example fermionic occupation constraints or positive terms that are discarded in the purely mean--field reduction.  The exact profile therefore provides a concrete reference state for future stability estimates around the radial obstruction.

The proof of \Cref{thm:main} has three main steps.  First, Newton's theorem and the logarithmic change of variables $t=\log r$ transform the radial minimization into a one--dimensional maximization problem for the kernel
\[
 k(t)=\ee^{-|t|}-\ee^{-2|t|}.
\]
Second, the Euler--Lagrange obstacle condition determines a compactly supported cosine profile.  Third, a conditional negative--definiteness estimate on the extremal support rules out all competing radial maximizers.  The last step is where the exact solution becomes global within the radial problem rather than merely stationary.

The finite--$Z$ step deliberately keeps the HPS finite--configuration and weighted kinetic inequalities unchanged.  The new input is the exact solution of the mean--field variational problem, together with a finite--$Z$ reoptimization that tracks the admissible interval for $N/Z$.  One small endpoint issue must be handled explicitly: for $Z\ge4$, Lieb's bound gives
\[
 \frac NZ<2+\frac1Z\le\frac94,
\]
which is sharper than the coarser interval used in the original numerical optimization and is sufficient to locate the relevant supremum at the left endpoint $N/Z=\beta_3^{-1}$.

The resulting extremal profile also clarifies a numerical feature already visible in HPS: radial trial densities close to $|x|^{-4}$ are nearly optimal.  In logarithmic radius, the exact extremizer is a compactly supported cosine perturbation of precisely that critical power law.

The paper is organized as follows.  \Cref{sec:radial} records the radial reduction.  \Cref{sec:log} introduces the logarithmic variational problem.  \Cref{sec:extremizer} constructs the candidate maximizer and proves its obstacle condition.  \Cref{sec:global} proves global optimality for the radial problem.  \Cref{sec:atomic} propagates the exact constant through the HPS finite--$N$ argument and proves \Cref{thm:finite-Z}.  \Cref{sec:outlook} discusses the significance and limitations of the result and the information required for improvements beyond the cubic mean--field obstruction.  The appendix isolates the technical points that are most sensitive in a formal proof audit.

\section{Radial reduction at $s=3$}\label{sec:radial}

We begin from the radial reduction proved in \cite{HPS2025}.  Let $\mu$ be radial and let $\nu$ denote its induced probability measure on radii, so that for every bounded radial test function $f$,
\[
 \int_{\R^3}f(|x|)\,\dd\mu(x)
 =\int_0^\infty f(r)\,\dd\nu(r).
\]
Newton's theorem gives
\[
 \int_{S^2}\int_{S^2}
 \frac{\dd\omega\,\dd\omega'}{(4\pi)^2|r\omega-s\omega'|}
 =\frac1{\max\{r,s\}}.
\]
Consequently the $s=3$ quotient becomes
\begin{equation}\label{eq:radial-quotient}
 \mathcal Q_3[\nu]
 :=
 \frac{\displaystyle
 \frac12\iint_{(0,\infty)^2}
 \frac{r^3+s^3}{\max\{r,s\}}\,\dd\nu(r)\dd\nu(s)}
 {\displaystyle\int_0^\infty r^2\,\dd\nu(r)}.
\end{equation}
Thus
\begin{equation}\label{eq:beta-radial}
 \beta_3=\inf_{\nu\in\mathcal A}\mathcal Q_3[\nu],
\end{equation}
where $\mathcal A$ is the radial image of the HPS admissible class.  The quotient is invariant under dilations: if $\nu_\lambda$ is the pushforward of $\nu$ under $r\mapsto\lambda r$, then
\[
 \mathcal Q_3[\nu_\lambda]=\mathcal Q_3[\nu].
\]

It is useful to isolate the kernel
\[
 K_3(r,s):=\frac{r^3+s^3}{2\max\{r,s\}}.
\]
For $0<r\le s$, writing $t=r/s$ gives
\[
 K_3(r,s)=\frac{s^2}{2}(1+t^3).
\]
The denominator contribution of the pair is naturally proportional to
\[
 \frac{r^2+s^2}{2}=\frac{s^2}{2}(1+t^2).
\]
The pointwise HPS estimate therefore uses
\[
 \frac{1+t^3}{1+t^2}\ge\min_{0\le t\le1}\frac{1+t^3}{1+t^2}.
\]
The minimum occurs at the unique $t_0\in(0,1)$ satisfying
\[
 t_0^3+3t_0-2=0,
\]
and yields
\[
 \frac{1+t_0^3}{1+t_0^2}=0.8941074569\ldots .
\]
Our purpose is to avoid this pairwise minimization.  A probability measure cannot, in general, arrange all of its pairs at the same worst radius ratio; the global measure constraint therefore contains additional information.  The logarithmic formulation below makes this constraint explicit.

\begin{remark}[Admissibility]\label{rem:admissibility}
HPS define $\beta_s$ over probability measures in $H^{-1}(\R^3)$ with the required finite moment.  The extremizer constructed below is absolutely continuous, supported on a compact annulus bounded away from the origin, and has bounded radial density there.  Hence it belongs to $H^{-1}(\R^3)$ and to the required moment class.  For the lower bound we may temporarily enlarge the radial admissible class to all positive radial probability measures for which \eqref{eq:radial-quotient} is finite.  Since the minimizer found in the enlarged class is HPS-admissible, equality of the two infima follows.
\end{remark}

\section{The logarithmic variational problem}\label{sec:log}

Set
\[
 t=\log r,
\]
and let $\bar\nu$ be the pushforward of $\nu$ under $r\mapsto t$.  Introduce the tilted positive measure $\sigma$ on $\R$ by
\begin{equation}\label{eq:sigma-def}
 \dd\sigma(t):=\ee^t\,\dd\bar\nu(t).
\end{equation}
Then
\[
 A[\sigma]:=\int_\R \ee^{-t}\,\dd\sigma(t)=1,
 \qquad
 B[\sigma]:=\int_\R \ee^t\,\dd\sigma(t)=\int_0^\infty r^2\,\dd\nu(r).
\]

For $t\ge u$, direct substitution gives
\[
 K_3(\ee^t,\ee^u)\,\ee^{-t}\ee^{-u}
 =\frac12\left(\ee^{t-u}+\ee^{-2(t-u)}\right).
\]
Thus, with $d=|t-u|$,
\[
 F(d):=\frac12(\ee^d+\ee^{-2d})
 =\cosh d-\frac12\bigl(\ee^{-d}-\ee^{-2d}\bigr).
\]
Define
\[
 k(d):=\ee^{-|d|}-\ee^{-2|d|}.
\]
Because
\[
 \iint \cosh(t-u)\,\dd\sigma(t)\dd\sigma(u)=A[\sigma]B[\sigma],
\]
we obtain
\[
 \mathcal Q_3[\nu]
 =1-\frac{1}{2}\frac{J[\sigma]}{A[\sigma]B[\sigma]},
\]
where
\[
 J[\sigma]:=\iint_{\R^2}k(t-u)\,\dd\sigma(t)\dd\sigma(u).
\]
Therefore
\begin{equation}\label{eq:M-def}
 \beta_3=1-\frac{M}{2},
 \qquad
 M:=\sup_{\sigma\ge0}\frac{J[\sigma]}{A[\sigma]B[\sigma]}.
\end{equation}
The quotient in \eqref{eq:M-def} is invariant under multiplication of $\sigma$ by a positive scalar and under translations $t\mapsto t+c$.  Hence we may impose the two normalizations
\begin{equation}\label{eq:normalization}
 A[\sigma]=B[\sigma]=1.
\end{equation}
In this gauge,
\begin{equation}\label{eq:M-normalized}
 M=\sup\left\{J[\sigma]:\sigma\ge0,
 \int\ee^{-t}\dd\sigma=\int\ee^t\dd\sigma=1\right\}.
\end{equation}
The supremum is strictly positive.  Indeed, for any $a>0$ the two--point measure
\[
 \sigma_a:=\frac{\delta_a+\delta_{-a}}{2\cosh a}
\]
satisfies $A[\sigma_a]=B[\sigma_a]=1$ and, since $k(0)=0$ and $k(2a)>0$,
\[
 J[\sigma_a]=\frac{k(2a)}{2\cosh^2 a}>0.
\]

\begin{lemma}[Existence of a maximizer]\label{lem:existence}
The supremum in \eqref{eq:M-normalized} is attained.
\end{lemma}

\begin{proof}
Let $(\sigma_n)$ be a maximizing sequence satisfying \eqref{eq:normalization}.  The moment constraints imply the uniform tail bounds
\[
 \sigma_n((R,\infty))\le \ee^{-R},
 \qquad
 \sigma_n(( -\infty,-R))\le \ee^{-R}.
\]
Moreover, by Cauchy--Schwarz,
\[
 \sigma_n(\R)^2
 \le\left(\int\ee^t\dd\sigma_n\right)
 \left(\int\ee^{-t}\dd\sigma_n\right)=1.
\]
Hence the sequence is tight and uniformly bounded in mass.  Passing to a subsequence, $\sigma_n\rightharpoonup\sigma$ narrowly as finite positive measures.  Narrow convergence together with convergence of the total masses implies
\[
 \sigma_n\otimes\sigma_n\rightharpoonup\sigma\otimes\sigma
\]
on $\R^2$.  Since $k(t-u)$ is bounded and continuous, it follows that
\[
 J[\sigma_n]\longrightarrow J[\sigma].
\]
In particular $J[\sigma]=M>0$, so $\sigma$ is nonzero.

Lower semicontinuity of the exponential moments gives $A[\sigma]\le1$ and $B[\sigma]\le1$.  Since $\sigma$ is a nonzero positive measure, both $A[\sigma]$ and $B[\sigma]$ are strictly positive.  If $A[\sigma]B[\sigma]<1$, first multiply $\sigma$ by $(A[\sigma]B[\sigma])^{-1/2}$ and then translate it by
\[
 c=\frac12\log\frac{A[\sigma]}{B[\sigma]}.
\]
The resulting positive measure $\widetilde\sigma$ satisfies $A[\widetilde\sigma]=B[\widetilde\sigma]=1$, while amplitude scaling and translation invariance give
\[
 J[\widetilde\sigma]
 =\frac{J[\sigma]}{A[\sigma]B[\sigma]}
 >J[\sigma]=M,
\]
a contradiction to the definition of $M$.  Hence $A[\sigma]B[\sigma]=1$.  Together with $A[\sigma]\le1$ and $B[\sigma]\le1$, this forces $A[\sigma]=B[\sigma]=1$.  Thus $\sigma$ is a normalized maximizer.
\end{proof}

The variational inequality can be obtained without invoking an abstract infinite-dimensional KKT theorem.

\begin{lemma}[Obstacle condition]
Let $\sigma$ be a normalized maximizer and let $M=J[\sigma]$.  Then
\begin{equation}\label{eq:obstacle}
 2(k*\sigma)(t)\le M(\ee^t+\ee^{-t})
 \qquad\text{for all }t\in\R,
\end{equation}
and equality holds on $\supp\sigma$.
\end{lemma}

\begin{proof}
Consider $\sigma+\varepsilon\delta_t$.  Since the quotient $J/(AB)$ is invariant under subsequent amplitude and translation renormalization, maximality gives
\[
 \frac{J[\sigma+\varepsilon\delta_t]}
 {A[\sigma+\varepsilon\delta_t]B[\sigma+\varepsilon\delta_t]}
 \le M.
\]
Because $k(0)=0$, differentiation at $\varepsilon=0+$ yields \eqref{eq:obstacle}.  Define
\[
 V(t):=M(\ee^t+\ee^{-t})-2(k*\sigma)(t)\ge0.
\]
Integrating against $\sigma$ gives
\[
 \int V\,\dd\sigma=M(A+B)-2J=0.
\]
Thus $V=0$ $\sigma$-almost everywhere, and continuity of $V$ implies equality on $\supp\sigma$.
\end{proof}

\section{Construction of the extremizer}\label{sec:extremizer}

Let
\begin{equation}\label{eq:a-star}
 a_*:=\frac1{\sqrt2}\arctan\sqrt2,
 \qquad q:=\sqrt2.
\end{equation}
Then
\begin{equation}\label{eq:tan-condition}
 \tan(q a_*)=q.
\end{equation}
Define
\begin{equation}\label{eq:M-star}
 M_*:=\frac{2}{1+3\ee^{2a_*}},
\end{equation}
and
\begin{equation}\label{eq:C-star}
 C_*:=\frac{3\sqrt3\,\ee^{a_*}}{1+3\ee^{2a_*}}.
\end{equation}
We consider the positive measure
\begin{equation}\label{eq:sigma-star}
 \dd\sigma_*(t)
 :=C_*\cos(qt)\,\mathbf 1_{[-a_*,a_*]}(t)\,\dd t.
\end{equation}
Since $qa_*\in(0,\pi/2)$, the density is strictly positive in the interior of its support.

\begin{proposition}[Normalization and contact equation]\label{prop:candidate}
The measure $\sigma_*$ satisfies
\[
 A[\sigma_*]=B[\sigma_*]=1,
\]
and
\begin{equation}\label{eq:contact-inside}
 (k*\sigma_*)(t)=M_*\cosh t,
 \qquad |t|\le a_*.
\end{equation}
Consequently
\[
 J[\sigma_*]=M_*.
\]
\end{proposition}

\begin{proof}
By symmetry it is enough to compute $B[\sigma_*]$.  Using $\tan(qa_*)=q$ and $q^2=2$,
\[
 \int_{-a_*}^{a_*}\ee^t\cos(qt)\,\dd t
 =\frac{3\ee^{a_*}+\ee^{-a_*}}{3\sqrt3}.
\]
The choice \eqref{eq:C-star} therefore gives $B[\sigma_*]=1$, and symmetry gives $A[\sigma_*]=1$.

For $\alpha>0$ define
\[
 Y_\alpha(t):=\int\ee^{-\alpha|t-u|}\,\dd\sigma_*(u).
\]
On the interior of the support,
\[
 (-\partial_t^2+\alpha^2)Y_\alpha
 =2\alpha C_*\cos(qt).
\]
Hence
\[
 Y_\alpha(t)
 =\frac{2\alpha C_*}{\alpha^2+q^2}\cos(qt)
 +A_\alpha\cosh(\alpha t),
 \qquad |t|<a_*,
\]
where evenness eliminates the $\sinh$ term.  For $\alpha=1$ and $\alpha=2$ the particular coefficients coincide:
\[
 \frac{2C_*}{1+q^2}
 =\frac{4C_*}{4+q^2}
 =\frac{2C_*}{3}.
\]
For $t\ge a_*$, $Y_\alpha$ is a pure multiple of $\ee^{-\alpha t}$, and therefore obeys the Robin condition
\[
 Y_\alpha'(a_*)=-\alpha Y_\alpha(a_*).
\]
For $\alpha=2$, condition \eqref{eq:tan-condition} makes the particular solution itself satisfy this Robin condition; hence $A_2=0$.  Subtracting $Y_2$ from $Y_1$ cancels the cosine particular solutions and gives
\[
 (k*\sigma_*)(t)=Y_1(t)-Y_2(t)=A_1\cosh t.
\]
Using the normalization $B[\sigma_*]=1$ in the exterior matching formula determines
\[
 A_1=\frac{2}{1+3\ee^{2a_*}}=M_*.
\]
This proves \eqref{eq:contact-inside}.  Finally,
\[
 J[\sigma_*]
 =\int(k*\sigma_*)\,\dd\sigma_*
 =M_*\int\cosh t\,\dd\sigma_*(t)
 =M_*,
\]
because $A=B=1$ implies $\int\cosh t\,\dd\sigma_*=1$.
\end{proof}

The contact equation extends to a strict obstacle inequality outside the support.

\begin{proposition}[Exact obstacle factorization]\label{prop:obstacle-star}
Let
\[
 V_*(t):=M_*(\ee^t+\ee^{-t})-2(k*\sigma_*)(t).
\]
Then $V_*(t)=0$ for $|t|\le a_*$ and $V_*(t)>0$ for $|t|>a_*$.  More precisely, for $t\ge a_*$ and
\[
 z:=\ee^{t-a_*}\ge1,
\]
one has
\begin{equation}\label{eq:factor-obstacle}
 \boxed{\displaystyle
 V_*(t)=M_*\ee^{a_*}\frac{(z-1)^2(z+2)}{z^2}.}
\end{equation}
The formula for $t\le-a_*$ follows by symmetry.
\end{proposition}

\begin{proof}
The vanishing on the support is exactly \eqref{eq:contact-inside}.  For $t\ge a_*$, all points $u$ in the support satisfy $u\le t$, so
\[
 (k*\sigma_*)(t)
 =\ee^{-t}\int\ee^u\,\dd\sigma_*(u)
 -\ee^{-2t}\int\ee^{2u}\,\dd\sigma_*(u).
\]
The first integral equals one.  The second is determined by the two contact conditions at $a_*$, namely $V_*(a_*)=V_*'(a_*)=0$.  Substituting the resulting coefficient and writing $z=\ee^{t-a_*}$ gives the exact factorization \eqref{eq:factor-obstacle}.  Its right side is nonnegative and vanishes only at $z=1$.
\end{proof}

Returning to the radial variable, \eqref{eq:sigma-def} gives
\begin{equation}\label{eq:nu-star}
 \dd\nu_*(r)
 =C_*r^{-2}\cos(\sqrt2\log r)\,
 \mathbf 1_{[\ee^{-a_*},\ee^{a_*}]}(r)\,\dd r.
\end{equation}
Thus the corresponding three--dimensional radial density is
\begin{equation}\label{eq:rho-star}
 \rho_*(x)
 =\frac{C_*}{4\pi}|x|^{-4}
 \cos(\sqrt2\log|x|)\,
 \mathbf 1_{\{\ee^{-a_*}\le|x|\le\ee^{a_*}\}}.
\end{equation}
This density is bounded on its annular support and hence belongs to the HPS admissible class.

\section{Global optimality}\label{sec:global}

We now prove that the candidate from \Cref{sec:extremizer} is the global maximizer of \eqref{eq:M-normalized}.

\subsection{Structure of an arbitrary maximizer}

Let $\sigma$ be a normalized maximizer, let
\[
 M:=J[\sigma]>0,
\]
and set $K:=k*\sigma$.

\begin{lemma}[Compact support and absence of atoms]\label{lem:compact-atomless}
Every normalized maximizer has compact support and is atomless.
\end{lemma}

\begin{proof}
The moment normalization implies
\[
 \sigma(\R)^2\le A[\sigma]B[\sigma]=1.
\]
Moreover $0\le k\le1/4$, hence
\[
 2K(t)\le\frac12\sigma(\R)\le\frac12.
\]
On the support, the obstacle condition \eqref{eq:obstacle} is an equality, so
\[
 M(\ee^t+\ee^{-t})=2K(t)\le\frac12.
\]
Since $M>0$, this confines the support to a compact interval.

Suppose that $\sigma$ has an atom of mass $m>0$ at $t_0$.  Since
\[
 k(x)=\ee^{-|x|}-\ee^{-2|x|}=|x|+O(x^2)
 \qquad (x\to0),
\]
the derivative of $K$ has the jump
\[
 K'(t_0+)-K'(t_0-)=2m.
\]
Thus the nonnegative obstacle
\[
 V(t)=M(\ee^t+\ee^{-t})-2K(t)
\]
has
\[
 V'(t_0+)-V'(t_0-)=-4m<0.
\]
But $t_0\in\supp\sigma$ implies $V(t_0)=0$, so $t_0$ is a minimum of the nonnegative function $V$; its one-sided derivatives must satisfy $V'(t_0-)\le0\le V'(t_0+)$, which is incompatible with the displayed negative jump.  Hence $\sigma$ has no atoms.
\end{proof}

Because $\sigma$ is atomless and compactly supported, convolution with the Lipschitz kernel $k$ is $C^1$.  Thus the obstacle is $C^1$ and the boundary of the support satisfies smooth pasting.

\begin{lemma}[Endpoint identities]\label{lem:endpoints}
Let
\[
 a:=\inf\supp\sigma,
 \qquad b:=\sup\supp\sigma.
\]
Then
\begin{equation}\label{eq:M-right}
 M=\frac{2}{1+3\ee^{2b}},
\end{equation}
and
\begin{equation}\label{eq:M-left}
 M=\frac{2\ee^{2a}}{\ee^{2a}+3}.
\end{equation}
Consequently $a+b=0$.
\end{lemma}

\begin{proof}
For $t\ge b$,
\[
 K(t)=\ee^{-t}\int\ee^u\,\dd\sigma(u)
 -\ee^{-2t}\int\ee^{2u}\,\dd\sigma(u)
 =\ee^{-t}-C_2\ee^{-2t},
\]
where $C_2:=\int\ee^{2u}\dd\sigma(u)$.  Since $b$ is a contact point and $V$ is $C^1$,
\[
 V(b)=V'(b)=0.
\]
Eliminating $C_2$ from these two equations yields \eqref{eq:M-right}.  The same argument at the left endpoint gives \eqref{eq:M-left}.  Equating the two expressions for $M$ gives
\[
 \ee^{2(a+b)}=1,
\]
hence $a+b=0$.
\end{proof}

We may therefore write the support of any maximizer as a compact subset of $[-c,c]$ with endpoints $\pm c$, and
\begin{equation}\label{eq:M-c}
 M=\frac{2}{1+3\ee^{2c}}.
\end{equation}

\subsection{Conditional negative definiteness on the extremal interval}

The Fourier transform of the kernel is
\begin{equation}\label{eq:khat}
 \widehat{k}(\xi)
 =\frac{2}{1+\xi^2}-\frac{4}{4+\xi^2}
 =\frac{2(2-\xi^2)}{(1+\xi^2)(4+\xi^2)}.
\end{equation}
Thus $k$ is not positive definite on the whole line.  What is needed is negativity on the codimension--two space of perturbations preserving the two exponential moments, restricted to the candidate support.

\begin{lemma}[Conditional negativity]\label{lem:conditional-negativity}
Let $\eta$ be a finite signed measure supported in $[-a_*,a_*]$ and satisfying
\begin{equation}\label{eq:eta-moments}
 \int\ee^t\,\dd\eta(t)=
 \int\ee^{-t}\,\dd\eta(t)=0.
\end{equation}
Then
\begin{equation}\label{eq:Jeta-negative}
 J[\eta]\le0,
\end{equation}
with equality only for $\eta=0$.
\end{lemma}

\begin{proof}
Define
\begin{equation}\label{eq:p-def}
 p(t):=\int\ee^{|t-u|}\,\dd\eta(u).
\end{equation}
The two moment conditions imply $p(t)=0$ for $|t|\ge a_*$.  Moreover $p\in H_0^1(-a_*,a_*)$ and, in the sense of distributions,
\begin{equation}\label{eq:p-eta}
 (\partial_t^2-1)p=2\eta.
\end{equation}
Using \eqref{eq:khat}, \eqref{eq:p-eta}, and Plancherel (first for smooth approximations and then by passage to the limit) gives the exact identity
\begin{equation}\label{eq:Jeta-identity}
 J[\eta]
 =\frac12\left[
 -\|p'\|_2^2+5\|p\|_2^2
 -18\left\langle p,(-\partial_t^2+4)^{-1}p\right\rangle
 \right].
\end{equation}
The resolvent term is nonnegative, hence
\[
 J[\eta]
 \le\frac12\bigl(5\|p\|_2^2-\|p'\|_2^2\bigr).
\]
Poincar\'e's inequality on an interval of length $2a_*$ yields
\[
 \|p'\|_2^2
 \ge\frac{\pi^2}{(2a_*)^2}\|p\|_2^2.
\]
Since
\begin{equation}\label{eq:length-threshold}
 2a_*=1.3510217177\ldots
 <\frac{\pi}{\sqrt5}=1.4049629462\ldots,
\end{equation}
we have $\pi^2/(2a_*)^2>5$, proving \eqref{eq:Jeta-negative}.  Equality forces $p=0$, and then \eqref{eq:p-eta} gives $\eta=0$.
\end{proof}

\subsection{Completion of the proof}

\begin{proof}[Proof of \Cref{thm:main}]
By \Cref{lem:existence}, choose a normalized global maximizer $\sigma$ with value $M$.  Since $\sigma_*$ is admissible and has value $M_*$,
\[
 M\ge M_*.
\]
Let the support endpoints of $\sigma$ be $\pm c$.  The endpoint identity \eqref{eq:M-c}, together with the strict monotonicity of $c\mapsto2/(1+3\ee^{2c})$, implies
\[
 c\le a_*.
\]
Hence both $\sigma$ and $\sigma_*$ are supported in $[-a_*,a_*]$.

Set
\[
 \eta:=\sigma-\sigma_*.
\]
Because both measures satisfy $A=B=1$, $\eta$ obeys \eqref{eq:eta-moments}.  Expanding the quadratic functional gives
\[
 J[\sigma]
 =J[\sigma_*]
 +2\int(k*\sigma_*)(t)\,\dd\eta(t)
 +J[\eta].
\]
On $[-a_*,a_*]$, \Cref{prop:candidate} gives $k*\sigma_*=M_*\cosh t$.  Therefore
\[
 2\int(k*\sigma_*)\,\dd\eta
 =M_*\int(\ee^t+\ee^{-t})\,\dd\eta=0.
\]
By \Cref{lem:conditional-negativity}, $J[\eta]\le0$.  Thus
\[
 M=J[\sigma]\le J[\sigma_*]=M_*.
\]
Together with $M\ge M_*$ this gives $M=M_*$ and, by the strict case in \Cref{lem:conditional-negativity}, $\sigma=\sigma_*$.  Hence the normalized maximizer is unique.

Finally, \eqref{eq:M-def} gives
\[
 \beta_3
 =1-\frac{M_*}{2}
 =1-\frac1{1+3\ee^{2a_*}}
 =\frac{3\ee^{2a_*}}{1+3\ee^{2a_*}},
\]
which is \eqref{eq:beta-exact}.  Undoing the logarithmic tilt yields \eqref{eq:nu-star}--\eqref{eq:rho-star}.  Translation in logarithmic radius corresponds to dilation in $r$, giving uniqueness up to scaling in the original problem.
\end{proof}

\section{Consequence for the excess--charge framework}\label{sec:atomic}

The role of $\beta_s$ in \cite{HPS2025} is to control the large--$N$ mean--field limit of the finite--configuration constants $\alpha_{N,s}$.  Their explicit $s=3$ estimate uses the pointwise radial lower bound
\[
 \beta_3\ge0.8941074569\ldots,
\]
whose reciprocal is
\[
 1.118433799\ldots .
\]
The exact evaluation in \Cref{thm:main} instead gives
\begin{equation}\label{eq:new-leading}
 \boxed{\displaystyle
 \beta_3^{-1}=1.086325175181735\ldots .}
\end{equation}
Thus, at the mean--field level, the coefficient available to the $s=3$ weighted method improves from approximately $1.1184$ to approximately $1.08633$.

This statement should be distinguished from a fully propagated finite--$Z$ atomic theorem.  HPS compare $\alpha_{N,3}$ with $\beta_3$ through quantitative smearing estimates and then combine that comparison with weighted kinetic--energy bounds.  Replacing the lower estimate for $\beta_3$ by its exact value improves the input to those arguments, but the optimal explicit lower--order term must be recomputed from the complete finite--$N$ inequalities.  We therefore record the following consequence in the form needed for such a recomputation.

\begin{corollary}[Exact mean--field input]\label{cor:hps-input}
Every finite--$N$ estimate in the HPS $s=3$ argument that depends monotonically on the mean--field constant $\beta_3$ may use the exact value
\[
 \beta_3=0.9205346822904167\ldots
\]
in place of the pointwise lower bound $0.8941074569\ldots$.
\end{corollary}

The improvement has a structural interpretation.  The pointwise estimate minimizes
\[
 g(t)=\frac{1+t^3}{1+t^2}
\]
for a single pair of radii.  It therefore permits every pair in the measure to realize the same worst ratio.  The exact optimizer shows that the true global compromise is a log--radial density
\[
 \rho_*(r)\propto r^{-4}\cos(\sqrt2\log r)
\]
on a finite annulus.  The resulting constant is strictly larger because the pairwise worst geometry cannot be realized simultaneously by a probability measure.

It is also useful to compare with the numerical trial family $\rho(r)\propto r^{-p}$ on a finite annulus considered in the HPS analysis.  The exact extremizer is centered on the critical power $p=4$ but acquires the log--oscillatory cosine factor forced by the Euler--Lagrange equation.  This explains why the simple $p\approx4$ family is already numerically close to the true value while remaining strictly nonoptimal.

\section{Discussion, significance, and scope}\label{sec:outlook}

The exact computation of $\beta_3$ and \Cref{thm:finite-Z} complete the finite--$Z$ propagation of the cubic HPS mean--field constant through the presently available finite--configuration and weighted kinetic estimates.  The two conclusions should be kept conceptually separate.  The first is an exact solution of the scalar radial mean--field variational problem.  The second is an explicit atomic bound obtained by inserting that exact constant into the existing HPS many--body framework and reoptimizing the finite--$Z$ bookkeeping.

At leading order this changes the coefficient from $1.1185$ to
\[
 \beta_3^{-1}=1.086325175\ldots .
\]
Measured against the conjectural coefficient $1$, the exact radial computation removes about $27\%$ of the gap left by the HPS leading coefficient.  This is a useful way to locate the contribution of the radial mean--field relaxation itself: no new kinetic inequality, localization estimate, or fermionic many--body input is used to obtain that reduction.  The remaining $Z^{1/3}$ correction is inherited from the finite--$N$ smearing/kinetic balance and is not eliminated by knowing $\beta_3$ exactly.

The explicit extremizer gives a second, qualitative conclusion.  The old pointwise estimate optimizes the interaction of a single pair of radii.  The true mean--field problem instead asks for one probability measure whose entire pair distribution is compatible at once.  The gap between the pointwise constant and the exact $\beta_3$ therefore measures a collective compatibility effect, and the minimizing distribution is an extended annular profile rather than a degenerate pair configuration.

In logarithmic radius the extremizer is a compactly supported cosine profile, while in physical space its density is proportional to
\[
 r^{-4}\cos(\sqrt2\log r)
\]
on a finite annulus.  This profile identifies the geometry responsible for the remaining scalar mean--field obstruction.  It also supplies a natural reference configuration for any future quantitative stability estimate: one can ask how much extra energy is forced when an admissible state departs from this log--radial profile.

The result also clarifies the limitation of the scalar approach.  The coefficient $\beta_3^{-1}>1$ is the value delivered by the cubic radial mean--field input together with the unchanged HPS finite--$N$ estimates.  We do not claim that it is optimal among all weighted multiplier arguments, and we do not prove the ionization conjecture $N_c(Z)\le Z+C$.  Any improvement beyond the coefficient obtained here within a strengthened cubic-weight framework must exploit information that disappears in the scalar radial reduction, for example genuinely many--fermion occupation constraints or additional positive terms retained from the weighted multiplier identity.

Thus the role of the present paper is twofold: it sharpens the best coefficient obtainable from the specific cubic mean--field input used in HPS, and it turns the remaining obstruction into an explicit extremal geometry rather than an unspecified constant loss.  That distinction is useful both for the quantitative excess--charge problem and for deciding where a subsequent improvement has to enter.

\appendix
\section{Proof audit and technical checkpoints}\label{app:audit}

This appendix records the points at which a formal variational or finite--$N$ calculation can fail and explains how they are handled in the proof above.  It is included deliberately: the exact numerical value of $\beta_3$ is meaningful only if the radial variational argument is genuinely global, and the improved atomic coefficient is meaningful only if the finite--$N$ propagation is numerically safe.

\subsection{Enlarging the HPS admissible class}

The HPS definition uses radial probability measures belonging to $H^{-1}(\R^3)$ and with finite second moment.  For the lower bound we enlarged the class after radial reduction.  This is harmless because the lower bound is proved for the larger class, while the explicit optimizer \eqref{eq:rho-star} is an $H^{-1}$ probability density with compact support bounded away from the origin.  Hence the lower and upper bounds meet inside the original HPS class.  HPS radial reduction identifies the value of the full variational infimum with the radial infimum; the uniqueness argument in the present manuscript, however, is only a uniqueness statement inside the radial class.

\subsection{Existence and loss of exponential moments}

Weak compactness alone does not preserve the constraints
\[
 A[\sigma]=B[\sigma]=1,
\]
because exponentially small masses can escape to large $|t|$ while carrying a nonvanishing exponential moment.  The proof of \Cref{lem:existence} therefore uses more than tightness.  First, a normalized two--point measure gives an explicit competitor with $J>0$, so the supremum $M$ is strictly positive.  Tightness and bounded total mass yield narrow convergence along a maximizing subsequence, and the corresponding product measures converge narrowly as well; because $k(t-u)$ is bounded and continuous, this gives $J[\sigma_n]\to J[\sigma]=M$.  If the weak limit had $A B<1$, amplitude and translation normalization would multiply $J$ by $(AB)^{-1}>1$, producing a normalized competitor with value larger than $M$.  This closes the only possible exponential-moment loss mechanism.

\subsection{The obstacle condition without an abstract KKT theorem}

The first variation is derived from the scale-- and translation--invariant quotient
\[
 \frac{J[\sigma]}{A[\sigma]B[\sigma]}.
\]
Adding an atom $\varepsilon\delta_t$ gives the global inequality directly.  Equality on the support then follows from the identity
\[
 \int\left[M(\ee^t+\ee^{-t})-2(k*\sigma)(t)\right]\dd\sigma(t)=0
\]
and nonnegativity of the bracket.  Thus no constraint qualification for an infinite-dimensional KKT theorem is required.

\subsection{Atoms and smooth pasting}

The cusp of $k$ at the origin is essential.  An atom of mass $m$ in $\sigma$ produces an upward derivative jump $2m$ in $k*\sigma$ and hence a downward jump $-4m$ in the derivative of the nonnegative obstacle.  This is incompatible with contact at a local minimum, excluding atoms.

For a compactly supported atomless finite measure, convolution with the Lipschitz kernel $k$ is $C^1$: the one-sided derivative discontinuity of $k'$ occurs only at $u=t$, which carries zero mass, and continuity follows by dominated convergence after splitting off a small neighborhood of $t$.  Therefore the obstacle is $C^1$, and at either extremal contact point one has both value and derivative equal to zero.  This justifies the smooth-pasting equations used in \Cref{lem:endpoints}.

\subsection{Conditional negativity for signed measures}

For a finite signed measure $\eta$ satisfying the two moment constraints, the function
\[
 p(t)=\int\ee^{|t-u|}\dd\eta(u)
\]
vanishes outside the candidate interval and belongs to $H_0^1$.  Distributionally,
\[
 (\partial_t^2-1)p=2\eta,
\]
so in particular $\eta\in H^{-1}(\R)$.  Fourier transformation gives
\[
 \widehat\eta(\xi)=-\frac{1+\xi^2}{2}\widehat p(\xi).
\]
Combining this identity with
\[
 \widehat{k}(\xi)=\frac{2(2-\xi^2)}{(1+\xi^2)(4+\xi^2)}
\]
gives the multiplier decomposition
\[
 \widehat{k}(\xi)|\widehat\eta(\xi)|^2
 =\frac12\left(-\xi^2+5-\frac{18}{\xi^2+4}\right)|\widehat p(\xi)|^2,
\]
which is exactly \eqref{eq:Jeta-identity}.  The finite-measure pairing can be justified by standard $H^{-1}$ approximation.  Thus the signed-measure step does not rely on a formal Fourier transform of a measure without a function-space interpretation.  The strict inequality ultimately uses only
\[
 2a_*<\frac{\pi}{\sqrt5}.
\]

\subsection{Why the global argument does not assume connected support}

No connectedness assumption on $\supp\sigma$ is made.  We define only its extremal points $a$ and $b$.  Compactness and smooth pasting at these two points yield the endpoint identities and hence $a+b=0$.  If a maximizer had value at least $M_*$, its extremal support points would lie inside $[-a_*,a_*]$, regardless of holes in the support.  The difference measure $\eta=\sigma-\sigma_*$ is therefore supported in the full candidate interval, exactly the hypothesis required in \Cref{lem:conditional-negativity}.

\subsection{Scope of the uniqueness statement}

The proof of \Cref{thm:main} shows that the normalized maximizer of the logarithmic radial problem is unique, and hence that the radial minimizer is unique up to dilation.  HPS radialization at $s=3$ proves equality between the full and radial infima, but the radialization inequality is not used here in a strict form.  Consequently the present paper does not claim that every minimizer in the full nonradial HPS admissible class must itself be radial.

\subsection{Why the exact $\beta_3$ may be inserted into the HPS finite--$N$ argument}

The finite--$Z$ theorem is not obtained by replacing a printed decimal inside the final HPS statement.  What is used is the symbolic cubic-weight chain in \cite[Section 7.3]{HPS2025}.  In particular, their finite--configuration estimate has the form
\[
 (1-r^2)N\beta_3-\frac1r
 \le
 \left(1+\frac{r^2}{5}+\frac{r^3}{35}\right)\alpha_{N,3}(N-1),
\]
and their kinetic estimate bounds the right--hand side by the same factor times
\[
 Z+c_0ZN^{-2/3}.
\]
Thus the true mean--field constant $\beta_3$ remains symbolic throughout the finite--$N$ argument.  The old pointwise lower bound is used only to turn $\beta_3^{-1}$ into the decimal $1.1185$ and to produce safe decimal lower--order coefficients.  Replacing that lower estimate by the exact value from \Cref{thm:main} is therefore logically legitimate.

\subsection{Endpoint check in the $Z^{1/3}$ optimization}

After the lower--order coefficients are removed, the HPS contradiction argument reduces to bounding
\[
 F_{\beta_3}(x)
 =
 3\left(\frac3{10}\right)^{1/3}\beta_3^{-2/3}x^{1/3}
 +c_0\beta_3^{-1}x^{-2/3},
 \qquad x=\frac NZ.
\]
Because
\[
 F'_{\beta_3}(x)
 =\frac{x^{-5/3}}{3}(Ax-2B)
\]
with $A,B>0$, the function has a unique interior minimum and no interior maximum.  Hence one must compare both endpoints of the admissible interval.

For the present theorem $Z\ge4$.  Lieb's bound gives
\[
 x<2+\frac1Z\le\frac94,
\]
so the relevant interval is $[\beta_3^{-1},9/4]$.  Directed evaluation of the explicit expressions gives
\[
 F_{\beta_3}(\beta_3^{-1})
 =3.81147042506\ldots,
 \qquad
 F_{\beta_3}(9/4)
 =3.78401155302\ldots.
\]
Thus the left endpoint is the true maximum on the interval used here, and the rounded coefficient $3.812$ is safe.

This sharper interval is important for the numerical audit.  If one unnecessarily enlarges the interval to $[\beta_3^{-1},5/2]$, the right endpoint is slightly larger for the exact value of $\beta_3$.  The $Z\ge4$ theorem avoids that artificial loss by retaining the sharper consequence of Lieb's estimate.

\subsection{Certified decimal margins}

All decimals in \Cref{thm:finite-Z} are chosen strictly on the safe side of explicit closed expressions:
\[
 \beta_3^{-1}=1.08632517518\ldots<1.086326,
\]
\[
 \bar a_2=0.01293244256\ldots<0.01294,\qquad
 \bar a_3=0.18184084198\ldots<0.18185,
\]
\[
 \bar a_4=0.01940243839\ldots<0.01941,
 \qquad
 \sup F_{\beta_3}=3.81147042506\ldots<3.812.
\]
These are evaluations of elementary explicit formulas involving $\exp$, $\arctan$, radicals and $\pi$; no numerical optimization enters the theorem.

\subsection{What is not claimed}

The paper proves a sharpened finite--$Z$ theorem within the existing HPS cubic-weight framework, but it does not prove the ionization conjecture $N_c(Z)\le Z+C$.  Nor does it claim that the coefficient $\beta_3^{-1}$ is optimal among all possible weighted multiplier arguments.  Any further reduction of the leading coefficient must use information not present in the scalar radial mean--field problem, for example a quantitative stability term coupled to fermionic occupation constraints.


\bibliographystyle{abbrv}
\bibliography{references}

\end{document}
